\chapter{ТУРШИЛТ БА ҮНЭЛГЭЭ}

\section{Үнэлгээний зорилго}

Энэ бүлэгт зан үйлд суурилсан дадал төлөвшүүлэх системийг бодит хэрэглэгчдээр туршуулсан үр дүнг нэгтгэн үнэлэв. Үнэлгээний зорилго нь системийн хэрэглээнд тохиромжтой байдал, дадал үүсгэх болон гүйцэтгэл бүртгэх урсгалын ойлгомжтой байдал, сануулгын хэрэгцээ, ахиц болон дадлын хүчний мэдээллийн ашиг тус, хэрэглэгчийн сэтгэл ханамжийг тодорхойлох явдал байв.

Мөн системийн цөм санаа болох өдөөгч нөхцөлд суурилсан дадал төлөвшүүлэлт, доод зорилт, төхөөрөмжийн мэдэгдэл, гүйцэтгэлийн бүртгэл, SRBAI-д тулгуурласан дадлын хүчний оноо, дасан зохицох зөвлөмж зэрэг боломжууд өдөр тутмын хэрэглээнд хэр ойлгомжтой, хэрэгтэй байгааг шалгахыг зорьсон.

\section{Судалгааны зохиомж}

Үнэлгээнд холимог арга зүйт, нэг бүлэгт суурилсан хэрэглэгчийн туршилтын зохиомж ашиглав. Оролцогчид системийг өөрийн сонгосон дадал дээр ашиглаж, туршилтын төгсгөлд Google Forms-оор судалгааны асуулгад хариулсан. Ингэснээр системийн дотоод бүртгэлийн өгөгдөл, хэрэглэгчийн өөрийн үнэлгээ, санал сэтгэгдлийг хамтатган авч үзэх боломж бүрдсэн.

Судалгаанд нийт 20 хэрэглэгч оролцсон. Оролцогчид системд бүртгүүлж, дор хаяж нэг дадал үүсгэн, хуваарь, цагийн хүрээ, өдөөгч нөхцөл, сануулга болон зорилтот хэмжилтийн мэдээллээ тохируулсан. Туршилтын явцад тэд гүйцэтгэлээ бүртгэх, сануулгад хариу өгөх, ахицын мэдээллээ харах, санал хүсэлт илгээх боломжтой байв.

\section{Оролцогчид}

Оролцогчдыг дадал төлөвшүүлэх хэрэгцээ өндөртэй зорилтот бүлгээс сонгосон. Үүнд их, дээд сургуулийн оюутан, залуу ажиллагч, өдөр тутмын хичээл, ажил, эрүүл мэнд, бүтээмжийн дадлаа тогтмол болгох сонирхолтой хэрэглэгчид багтсан. Оролцогч бүр өөрт хэрэгтэй гэж үзсэн дадлыг сонгосон тул бүртгэгдсэн дадлууд нь дасгал хөдөлгөөн, ус уух, унших, хичээл хийх, эрт унтах, хувийн бүтээмжийн жижиг үйлдэл зэрэг төрөл бүрийн чиглэлтэй байв.

\section{Туршилтын явц}

Туршилтыг дараах үндсэн үе шаттайгаар зохион байгуулсан.

\begin{enumerate}
    \item \textbf{Танилцуулга ба бүртгэл:} Оролцогчдод системийн зорилго, үндсэн боломжуудыг тайлбарлаж, бүртгэл болон нэвтрэх урсгалыг ашиглуулсан.
    \item \textbf{Дадал үүсгэх:} Хэрэглэгч бүр дадлын нэр, хэмжилтийн төрөл, доод зорилт, хуваарь, өдөөгч нөхцөл, цагийн хүрээ, сануулгын тохиргоог бөглөсөн.
    \item \textbf{Өдөр тутмын хэрэглээ:} Оролцогчид дадлаа гүйцэтгэх үедээ системд бүртгэл хийж, сануулга ирсэн тохиолдолд ``хийсэн'' эсвэл ``хойшлуулах'' үйлдэл ашигласан.
    \item \textbf{Ахиц ба зөвлөмж харах:} Хэрэглэгчид үндсэн, дадлын дэлгэрэнгүй, ахицын шинжилгээ, профайл дэлгэцүүдээр дамжуулан ахиц, дараалал, дадлын хүч, амжилт, зөвлөмжийн мэдээлэлтэй харилцсан.
    \item \textbf{Эцсийн судалгаа:} Туршилтын дараа хэрэглэгчид профайл хэсэгт байршуулсан Google Forms холбоосоор ашиглахад хялбар байдал, үр ашиг, сануулгын тохиромжтой байдал, сайжруулах саналын талаар хариулт өгсөн.
\end{enumerate}

\section{Цуглуулсан өгөгдөл}

Үнэлгээнд ашигласан өгөгдлийг хоёр үндсэн бүлэгт хуваав. Нэгдүгээрт, системийн ажиллагаанаас автоматаар үүсэх тоон өгөгдөл; хоёрдугаарт, хэрэглэгчийн өөрийн үнэлгээ болон чөлөөт санал сэтгэгдэл юм.

\begin{table}[htbp]
\centering
\small
\setlength{\tabcolsep}{4pt}
\begin{tabular}{|p{4.0cm}|p{8.8cm}|}
\hline
\textbf{Өгөгдлийн төрөл} & \textbf{Тайлбар} \\ \hline
Гүйцэтгэлийн бүртгэл & Дадал хийсэн эсэх, бодит утга, бүртгэсэн хугацаа, доод зорилт болон үндсэн зорилтын биелэлт \\ \hline
Сануулгын өгөгдөл & Илгээсэн сануулга, хийсэн/хойшлуулах үйлдэл, сануулгын хугацаа, холбогдсон сануулгын мэдээлэл \\ \hline
Өдөөгч нөхцөл ба нөхцөл байдлын өгөгдөл & Цагийн хүрээ, байршил, өмнөх үйлдэл, тухайн нөхцөлд гүйцэтгэсэн эсэх \\ \hline
Дадлын хүчний өгөгдөл & SRBAI үнэлгээ, тогтвортой байдал, нөхцөл байдлын тогтвортой байдал, нийлмэл дадлын хүчний оноо \\ \hline
Оролцооны өгөгдөл & XP, амжилтын тэмдэг, дараалал, профайл болон ахицын шинжилгээний дэлгэцийн хэрэглээ \\ \hline
Судалгааны өгөгдөл & Ашиглахад хялбар байдал, сэтгэл ханамж, сануулгын үр нөлөө, сайжруулах санал \\ \hline
\end{tabular}
\caption{Үнэлгээнд ашигласан өгөгдлийн төрлүүд}
\label{tab:evaluation-data-types}
\end{table}

\section{Үнэлгээний шалгуур үзүүлэлтүүд}

Системийг дараах шалгуур үзүүлэлтээр үнэлэв.

\begin{itemize}
    \item \textbf{Хэрэглээний ойлгомжтой байдал:} бүртгэл, дадал үүсгэх, гүйцэтгэл бүртгэх урсгал хэрэглэгчид ойлгомжтой эсэх;
    \item \textbf{Дадал төлөвшүүлэх дэмжлэг:} өдөөгч нөхцөл, доод зорилт, сануулга, шалтгаан, алхам зэрэг тохиргоо нь дадлаа хэрэгжүүлэхэд тусалсан эсэх;
    \item \textbf{Сануулгын тохиромжтой байдал:} төхөөрөмжийн мэдэгдэл зөв цагт, хэрэгтэй мэдээлэлтэй ирсэн эсэх;
    \item \textbf{Ахицын мэдээллийн хэрэгцээ:} үндсэн, ахицын шинжилгээ, дадлын дэлгэрэнгүй дэлгэцүүдийн мэдээлэл хэрэглэгчид өөрийн явцыг ойлгоход тусалсан эсэх;
    \item \textbf{Урамшууллын нөлөө:} XP, амжилтын тэмдэг, дараалал, амжилтын мэдээлэл нь хэрэглэгчийн үргэлжлүүлэх сэдэлд эерэг нөлөөтэй эсэх;
    \item \textbf{Системийн ерөнхий сэтгэл ханамж:} хэрэглэгчид цаашид үргэлжлүүлэн ашиглах хүсэлтэй эсэх.
\end{itemize}

\section{Судалгааны үр дүн}

\subsection{Ашиглахад хялбар байдал}

Судалгааны хариултаас харахад хэрэглэгчид бүртгүүлэх, дадал үүсгэх, тухайн өдрийн дадлаа харах, гүйцэтгэл бүртгэх урсгалыг ерөнхийдөө ойлгомжтой гэж үнэлсэн. Ялангуяа үндсэн дэлгэц дээр өдрийн дадлууд нэг дор харагдах, дадлын картаас шууд гүйцэтгэл бүртгэх боломж нь өдөр тутмын хэрэглээнд тохиромжтой гэж үзэгдсэн.

Дадал үүсгэх маягт нь олон тохиргоотой боловч өгүүлбэрэн бүтэцтэй зохион байгуулагдсан. Энэ нь хэрэглэгчид ``юу хийх'', ``хэзээ хийх'', ``ямар нөхцөлд хийх'', ``яагаад хийх'' гэсэн асуултаар дадлаа тодорхойлоход тусалсан. Гэсэн хэдий ч зарим хэрэглэгч анх удаа ашиглах үед нэмэлт заавар хэрэгтэй гэж санал өгсөн. Иймээс эхний танилцуулга болон товч тусламжийн тайлбарыг цаашид нэмэгдүүлэх шаардлагатай.

\subsection{Өдөөгч нөхцөл болон сануулгын үр нөлөө}

Оролцогчдын саналд цагийн хүрээ, өмнөх үйлдэл, төхөөрөмжийн мэдэгдэл зэрэг боломжууд дадлаа марталгүй хийхэд тусалсан гэж давтагдан дурдагдсан. Ялангуяа ``өглөө'', ``хичээлийн дараа'', ``унтахын өмнө'' гэх мэт өдөр тутмын тогтсон нөхцөлтэй холбож дадал үүсгэх нь бодит амьдралын хуваарьтайгаа уялдуулахад дөхөм байв.

Сануулгын хувьд ``хийсэн'' болон ``хойшлуулах'' гэсэн хоёр үйлдэл хангалттай энгийн байсан. Хийсэн үйлдэл нь сануулгатай холбоотой гүйцэтгэлийг автоматаар бүртгэх тул хэрэглэгчээс нэмэлт алхам шаардахгүй. Харин хойшлуулах үйлдэл нь тухайн үед хийх боломжгүй боловч бүр мөсөн мартахгүй байх хэрэгцээг хангаж байна. Үүнээс үзэхэд сануулгыг зөвхөн мэдэгдэл илгээх бус, хэрэглэгчийн нөхцөл байдалд тохирсон богино үйлдэлтэй холбох нь хэрэглээний үр ашгийг нэмэгдүүлсэн.

\subsection{Ахицын шинжилгээ ба дадлын хүчний мэдээлэл}

Үндсэн болон дадлын дэлгэрэнгүй дэлгэцүүд дээр дараалал, хуанли, гүйцэтгэлийн түүх, цагийн хүрээ, өдөөгч нөхцөл, шалтгаан зэрэг мэдээллийг нэг дор харуулсан нь хэрэглэгчид өөрийн дадлын явцыг ойлгоход тусалсан. Ахицын шинжилгээний хуанли болон идэвхийн дүрслэл нь аль өдрүүдэд тогтвортой байсан, аль үед тасалдсан гэдгийг харахад хэрэгтэй гэж үнэлэгдсэн.

Дадлын хүчний оноо нь SRBAI, гүйцэтгэлийн тогтвортой байдал, нөхцөл байдлын тогтвортой байдал гэсэн гурван бүрэлдэхүүнтэй. Тиймээс зөвхөн ``хийсэн тоо''-гоор бус, дадал автоматжиж байгаа эсэхийг илүү өргөн өнцгөөс тайлбарлах боломж өгсөн. Хэрэглэгчийн хувьд энэ оноо нь өөрийн ахицыг нэг тоон үзүүлэлтээр харах давуу талтай боловч бүрэлдэхүүн хэсгүүдийг илүү энгийн хэлээр тайлбарлах хэрэгцээ ажиглагдсан.

\subsection{Урамшуулал, XP ба амжилтын тэмдэг}

Профайл хэсэгт XP, амжилтын тэмдэг, амжилтын мэдээллийг харуулсан нь хэрэглэгчийн хийсэн үйлдлийг баталгаажуулж, үргэлжлүүлэх сэдлийг нэмэгдүүлэхэд эерэг нөлөөтэй байв. Анхны гүйцэтгэл, дараалсан өдрийн дараалал, нийт гүйцэтгэлийн үе шат зэрэг амжилтын тэмдэг нь богино хугацааны туршилтад ч хэрэглэгчид ахиц мэдрүүлэх боломжтой байсныг харуулсан.

Гэхдээ урамшууллын систем нь хэт төвөгтэй эсвэл гол дадлын зорилгоос анхаарал сарниулахгүй байх шаардлагатай. Иймээс XP болон амжилтын тэмдгийг хэрэглэгчийг шахах бус, ахицыг эерэгээр бататгах нэмэлт элемент байдлаар хадгалах нь тохиромжтой гэж үзэв.

\subsection{Санал хүсэлт ба сайжруулах хэрэгцээ}

Судалгааны чөлөөт саналуудаас дараах нийтлэг сайжруулалтын чиглэлүүд тодорсон.

\begin{itemize}
    \item дадал үүсгэх үед алхам бүрийн богино тайлбар, жишээ нэмэх;
    \item дадлын хүчний оноо болон зөвлөмжийн тайлбарыг илүү энгийн хэлбэрээр харуулах;
    \item мэдэгдлийн давтамж, цагийг хэрэглэгч илүү нарийвчлан тохируулах боломжийг нэмэгдүүлэх;
    \item ахицын шинжилгээний хэсэгт долоо хоног, сар, нийт хугацааны харьцуулсан харагдацыг илүү тодорхой болгох;
    \item профайл хэсгийн амжилт, судалгаа, санал хүсэлтийн хэсгийг илүү товч, эмх цэгцтэй байлгах.
\end{itemize}

\section{Хэлэлцүүлэг}

Туршилтын үр дүнгээс харахад зан үйлд суурилсан дадал төлөвшүүлэх систем нь хэрэглэгчид дадлаа тодорхойлох, тохиромжтой өдөөгч нөхцөл сонгох, өдрийн гүйцэтгэлээ бүртгэх, ахицаа харах үндсэн хэрэгцээг хангаж байна. Системийн онцлог нь ердийн хийх ажлын жагсаалт эсвэл дадал хянагчаас ялгаатайгаар дадлыг зөвхөн жагсаалт хэлбэрээр бус, цаг, байршил, өмнөх үйлдэл, доод зорилт, шалтгаан, сануулга, дадлын хүчний оноо зэрэг зан үйлийн бүрэлдэхүүн хэсгүүдтэй холбож загварчилсанд оршино.

Сануулгын хэсэг хэрэглэгчийг дадлаа мартахгүй байхад тусалж байгаа боловч урт хугацаанд сануулгаас хэт хамаарах эрсдэлтэй. Иймээс системд дадлын хүч нэмэгдэхийн хэрээр сануулгын эрчмийг бууруулах бодлого хэрэгжүүлсэн нь онолын хувьд зөв чиглэл юм. Богино хугацааны хэрэглэгчийн туршилтаар энэ боломжийн урт хугацааны нөлөөг бүрэн хэмжих боломж хязгаарлагдмал байсан ч системийн архитектур ирээдүйн урт хугацааны судалгаанд ашиглахад бэлэн болсон.

Мөн хэрэглэгчийн саналуудаас харахад системийн мэдээлэл их байх тусам тайлбар, зааварчилгаа чухал болдог. Тиймээс дараагийн хувилбарт эхний танилцуулга, тухайн үзүүлэлтийн товч тайлбар, онооны тайлбар, зөвлөмжийг хэрэгжүүлэх алхам болгон хувиргах зэрэг хэрэглэгчийн туршлагын сайжруулалт хийх шаардлагатай.

\section{Хязгаарлалт}

Энэхүү үнэлгээ нь 20 хэрэглэгчтэй богино хугацааны хэрэглэгчийн туршилтад тулгуурласан. Иймээс үр дүнг бүх хэрэглэгчийн бүлэгт шууд ерөнхийлөх боломж хязгаарлагдмал. Мөн оролцогчдын сонгосон дадал, хэрэглээний орчин, өдрийн хуваарь ялгаатай тул дадал бүрийн амжилтын түвшинг шууд харьцуулахад болгоомжтой хандах шаардлагатай.

Судалгааны өгөгдөл нь системийн бүртгэл болон хэрэглэгчийн өөрийн үнэлгээнд тулгуурласан тул өөрийн үнэлгээний хазайлт үүсэх боломжтой. Түүнчлэн дадал автоматжилт нь олон долоо хоног, зарим тохиолдолд олон сар шаардах үйл явц тул SRBAI болон нийлмэл дадлын хүчний онооны урт хугацааны өөрчлөлтийг бүрэн батлахад илүү урт хугацааны судалгаа хэрэгтэй.

\section{Бүлгийн дүгнэлт}

20 хэрэглэгчийн оролцоотой туршилтын үр дүнд системийн үндсэн хэрэгжүүлэлт бодит хэрэглээнд ашиглагдах боломжтой болох нь харагдав. Хэрэглэгчид дадал үүсгэх, өдөр тутмын гүйцэтгэл бүртгэх, сануулга авах, ахицаа харах, амжилтын тэмдгээр урамших зэрэг гол урсгалуудыг ашигласан. Судалгааны саналуудаас харахад системийн хамгийн хэрэгтэй хэсгүүд нь өдрийн дадлын үндсэн дэлгэц, төхөөрөмжийн мэдэгдэл, гүйцэтгэлийн бүртгэл, ахицын хуанли ба шинжилгээ, энгийн урамшууллын элементүүд байв.

Үүний зэрэгцээ дадал үүсгэх урсгалын тайлбар, дадлын хүчний онооны ойлгомжтой байдал, мэдэгдлийн нарийвчилсан тохиргоо, ахицын шинжилгээний харьцуулсан дүрслэл зэрэг сайжруулах чиглэлүүд тодорхой болсон. Иймээс энэхүү систем нь дипломын ажлын хүрээнд тавьсан зорилгыг хангаж, зан үйлд суурилсан дадал төлөвшүүлэх програм хангамжийн бүрэн ажиллагаатай хувилбар болон хэрэгжсэн гэж дүгнэж байна.


